\NativeSlot{45}{79}{%
\ExamGroup{二、选择（将正确答案的序号填在括号里）。{\normalfont\kaishu\fontsize{10.8bp}{14bp}\selectfont （每题3分，共18分）}}

% p086_s01

\ExamQuestion{1.}{5}%
将图形\raisebox{-11bp}{{\setlength{\NativeFigureHeight}{27bp}\NativeImage{figures/p086_s01-01.tex}}}\,顺时针旋转\(90^\circ\)后可以得到下面图形\mbox{（\hspace{1.8em}）}。\par
\par\noindent%
\ExamTile{.242}{%
A. \raisebox{-14bp}{{\setlength{\NativeFigureHeight}{39bp}\NativeImage{figures/p086_s01-02.tex}}}
}\hfill%
\ExamTile{.242}{%
B. \raisebox{-14bp}{{\setlength{\NativeFigureHeight}{39bp}\NativeImage{figures/p086_s01-03.tex}}}
}\hfill%
\ExamTile{.242}{%
C. \raisebox{-14bp}{{\setlength{\NativeFigureHeight}{39bp}\NativeImage{figures/p086_s01-04.tex}}}
}\hfill%
\ExamTile{.242}{%
D. 无法确定
}\par
}
\NativeSlot{179}{61}{%
% p086_s02

\ExamQuestion{2.}{5}%
在同一平面内的\(a,b,c\)三条直线，直线\(a\)与直线\(b\)都和直线\(c\)相交成直角，则直线\(a\)与直线\(b\)\mbox{（\hspace{1.8em}）}。\par
\ExamOptions{A. 互相平行}{B. 互相垂直}{C. 不平行也不垂直}{D. 无法确定}
}
\NativeSlot{282}{40}{%
% p086_s03

\ExamQuestion{3.}{5}%
等腰三角形中，有一个内角是\(40^\circ\)，另外两个内角是\mbox{（\hspace{1.8em}）}。\par
\ExamOptions{A. \(40^\circ\)和\(100^\circ\)}{B. \(70^\circ\)和\(70^\circ\)}{C. \(40^\circ\)和\(100^\circ\)或者\(70^\circ\)和\(70^\circ\)}{D. 无法确定}
}
\NativeSlot{357}{43}{%
% p086_s04

\ExamQuestion{4.}{5}%
把一个长方形框架拉成一个平行四边形，平行四边形的周长与原长方形的周长\mbox{（\hspace{1.8em}）}。\par
\ExamOptions{A. 大}{B. 小}{C. 一样大}{D. 无法比较大小}
}
\NativeSlot{434}{47}{%
% p086_s05

\ExamQuestion{5.}{5}%
下列图形中，不是轴对称图形的是\mbox{（\hspace{1.8em}）}。\par
\BookFigure{figures/p086_s05-01.tex}
}
\NativeSlot{516}{57}{%
% p086_s06

\ExamQuestion{6.}{5}%
以下是某手机地图APP中的四个工具标识图，哪个不是轴对称图形？\mbox{（\hspace{1.8em}）}\par
\BookFigure{figures/p086_s06-01.tex}
}
\NativeSlot{615}{107}{%
\ExamGroup{三、判断（对的打“√”，错的打“×”）。{\normalfont\kaishu\fontsize{10.8bp}{14bp}\selectfont （共12分）}}

% p086_s07

\ExamQuestion{1.}{5}%
过直线外一点只能画一条与已知直线互相平行的直线。\mbox{（\hspace{1.8em}）}\par

% p086_s08

\ExamQuestion{2.}{5}%
不相交的两条直线互相平行。\mbox{（\hspace{1.8em}）}\par

% p086_s09

\ExamQuestion{3.}{5}%
一个等腰三角形，底角是\(25^\circ\)，它一定是钝角三角形。\mbox{（\hspace{1.8em}）}\par

% p086_s10

\ExamQuestion{4.}{5}%
旋转后图形的大小、形状都改变了。\mbox{（\hspace{1.8em}）}\par

% p086_s11

\ExamQuestion{5.}{5}%
相交后的四个角都是直角的两条直线一定互相垂直。\mbox{（\hspace{1.8em}）}\par

% p086_s12

\ExamQuestion{6.}{5}%
两个等底等高的梯形一定能拼成一个平行四边形。\mbox{（\hspace{1.8em}）}\par
}
